
\medskip

%Un professeur propose un jeu à ses élèves.
%
%Ils doivent tirer un jeton dans une boîte de leur choix et gagnent lorsqu'ils tombent sur un jeton noir. 
%
%Le professeur leur précise que :
%
%\setlength\parindent{1cm}
%\begin{itemize}
%\item La boîte A contient la jetons dont 1 jeton noir ;
%\item La boîte B contient 15\,\% de jetons noirs ;
%\item La boîte C contient exactement $350$ jetons blancs et $50$ jetons noirs.
%\end{itemize}
%\setlength\parindent{0cm}
%
%\medskip
%
%Les jetons sont indiscernables au toucher. Une fois que l'élève a choisi sa boîte, le tirage se fait au hasard.
%
%\medskip

\begin{enumerate}
\item %Montrer que, dans la boîte C, la probabilité de tirer un jeton noir est $\dfrac{1}{8}$.
La probabilité de tirer un jeton noir dans la boîte C est égale à $\dfrac{50}{350 + 50} = \dfrac{50}{400}= \dfrac{50 \times 1}{50 \times 8} = \dfrac{1}{8}$.
\item %C'est le tour de Maxime. Dans quelle boîte a-t-il intérêt à tenter sa chance ? Justifier la réponse.
La probabilité de tirer un jeton noir dans la boîte A est égale à $\dfrac{1}{10} = 0,1$ ;

la probabilité de tirer un jeton noir dans la boîte B est égale à $\dfrac{15}{100} = 0,15$ et 

La probabilité de tirer un jeton noir dans la boîte C est égale à $\dfrac{1}{8} = 0,125$.

Comme $0,1 < 0,125 < 0,15$, Maxime a intérêt à choisir la boîte B.
\item %La boîte B contient $18$ jetons noirs. Combien y a-t-il de jetons au total dans cette boîte  ?
On a pour $n$ jetons en tout : $0,15 = \dfrac{15}{n}$ soit $0,15n = 18$ ou $n = \dfrac{18}{0,15} = 120$.

Il y a 120 jetons dans la boîte B dont 18 noirs. 
\item  %On ajoute $10$ jetons noirs dans la boîte C. Déterminer le nombre de jetons blancs à ajouter dans la boîte C pour que la probabilité de tirer un jeton noir reste égale à $\dfrac{1}{8}$.
Si on ajoute $b$ jetons blancs dans la boîte C, on a donc :

$\dfrac{50 + 10}{350 + 50 + 10 + b} = \dfrac{1}{8}$ ou $\dfrac{60}{410 + b} = \dfrac{1}{8}$, d'où on déduit : $8 \times 60 = 410 + b$ ou $480 = 410 + b$ et 

$b = 480 - 410 = 70$. Il faut ajouter 70 jetons blancs.
\end{enumerate}

\bigskip

